%% An Introduction to LaTeX Thesis Template of Yangtze University
%%
%% Created by gis-xh

% 导言区

% 使用ctexart类，12pt字号，a4纸张，单页模式
\documentclass[UTF8, 12pt, a4paper, oneside]{ctexart}
% 使用inputenc宏包，设置文档编码方式
% \usepackage[UTF8]{inputenc}
\usepackage{amsmath, amsthm, amssymb, graphicx}
\usepackage[bookmarks=true, colorlinks, citecolor=blue, linkcolor=black]{hyperref}

% 使用geometry宏包，设置页边距
\usepackage{geometry}
\geometry{left=3.0cm,right=3.0cm,top=3.0cm,bottom=2.5cm}
% 设置行间距
\linespread{1.25}
% 设置标准字间距
\usepackage{microtype}

% 设置页眉页脚
\usepackage{fancyhdr} % 引入宏包
\pagestyle{fancy} % 设置页面样式为 fancy
\fancyhf{} % 清除默认的页眉页脚
% 页眉中间显示小五号宋体的章节名称
\fancyhead[C]{\zihao{-5}\songti\leftmark}
% 页脚中间显示小五号宋体的页码
\fancyfoot[C]{\zihao{-5}\songti\thepage}

\usepackage{titlesec} 
% 定义封面内容
% \title{A Sample LaTeX Document} \author{测绘QX221 谢泓} \date{\today}
% % 定义封面样式
% \titleformat{\title}{\centering\Huge\bfseries}{\thetitle}{1em}{}
% \titleformat{\author}{\centering\Large}{\theauthor}{1em}{}
% \titleformat{\date}{\centering}{\thedate}{1em}{}

\begin{document}

%第一个封面，没有图片logo
\begin{titlepage}
    \centering
    \vspace*{5cm} %设置标题距离页面顶部的距离
    \includegraphics[height=1.69cm, width=6.68cm]{./figures/Yangtze_University.png}
    
    \vspace{1cm}
    {\Huge 硕 士 研 究 生 学 位 论 文}\par %设置标题字体大小和换行
    \vspace{1cm} %设置副标题距离标题的距离
    {\Large 基于 Latex 的模板}\par %设置副标题字体大小和换行
    \vfill %填充垂直空白
    {\Large 作者：我}\par %设置作者信息字体大小和换行
    \vspace{0.5cm} %设置日期距离作者信息的距离
    {\Large 论文起止日期：2022年1月1日至}\par %设置日期字体大小和换行
\end{titlepage}

%第二个封面，有图片logo
\begin{titlepage}
    \centering
    \vspace*{3cm} %设置图片logo距离页面顶部的距离
    % 插入图片logo
    \includegraphics[height=2.86cm,width=2.86cm]{./figures/Yangtze_University_logo.png}
    \par
    \includegraphics[height=1.69cm, width=6.68cm]{./figures/Yangtze_University.png}
    
    \vspace{1cm} %设置标题距离图片logo的距离
    {\Huge 这是我的标题}\par %设置标题字体大小和换行
    \vspace{1cm} %设置副标题距离标题的距离
    {\Large 这是我的副标题}\par %设置副标题字体大小和换行
    \vfill %填充垂直空白
    {\Large 作者：我}\par %设置作者信息字体大小和换行
    \vspace{0.5cm} %设置日期距离作者信息的距离
    {\Large 日期：2022年1月1日}\par %设置日期字体大小和换行
\end{titlepage}

% 中文摘要

论文的摘要是对论文研究内容和成果的高度概括。
摘要应对论文所研究的问题及其研究目的进行描述，对研究方法和过程进行简单介绍，对研究成果和所得结论进行概括。
摘要应具有独立性和自明性，其内容应包含与论文全文同等量的主要信息。
使读者即使不阅读全文，通过摘要就能了解论文的总体内容和主要成果。

论文摘要的书写应力求精确、简明。
切忌写成对论文书写内容进行提要的形式，尤其要避免“第 1 章……；第 2 章……；……”这种或类似的陈述方式。

关键词是为了文献标引工作、用以表示全文主要内容信息的单词或术语。
关键词不超过 5 个，每个关键词中间用分号分隔。
 % 引入中文摘要
\newpage % 另起一页

\tableofcontents % 生成目录
\newpage % 另起一页

% 英文摘要

An abstract of a dissertation is a summary and extraction of research work and contributions. Included in an abstract should be description of research topic and research objective, brief introduction to methodology and research process, and summary of conclusion and contributions of the research. An abstract should be characterized by independence and clarity and carry identical information with the dissertation. It should be such that the general idea and major contributions of the dissertation are conveyed without reading the dissertation.

An abstract should be concise and to the point. It is a misunderstanding to make an abstract an outline of the dissertation and words ``the first chapter'', ``the second chapter'' and the like should be avoided in the abstract.

Keywords are terms used in a dissertation for indexing, reflecting core information of the dissertation. An abstract may contain a maximum of 5 keywords, with semi-colons used in between to separate one another. % 引入英文摘要
\newpage % 另起一页

% Chapter 1

\chapter{论文主要部分的写法}
研究生学位论文撰写，除表达形式上需要符合一定的格式要求外，内容方面上也要遵循一些共性原则。

通常研究生学位论文只能有一个主题（不能是几块工作拼凑在一起），该主题应针对某学科领域中的一个具体问题展开深入、系统的研究，并得出有价值的研究结论。学位论文的研究主题切忌过大，例如，“中国国有企业改制问题研究”这样的研究主题过大，因为“国企改制”涉及的问题范围太广，很难在一本研究生学位论文中完全研究透彻。

\section{论文的语言及表述}
除国际研究生外，学位论文一律须用汉语书写。学位论文应当用规范汉字进行撰写，除古汉语研究中涉及的古文字和参考文献中引用的外文文献之外，均采用简体汉字撰写。

国际研究生一般应以中文或英文书写学位论文，格式要求同上。论文须用中文封面。

研究生学位论文是学术作品，因此其表述要严谨简明，重点突出，专业常识应简写或不写，做到立论正确、数据可靠、说明透彻、推理严谨、文字凝练、层次分明，避免使用文学性质的或带感情色彩的非学术性语言。

论文中如出现一个非通用性的新名词、新术语或新概念，需随即解释清楚。

\section{论文题目的写法}
论文题目应简明扼要地反映论文工作的主要内容，力求精炼、准确，切忌笼统。论文题目是对研究对象的准确、具体描述，一般要在一定程度上体现研究结论，因此，论文题目不仅应告诉读者这本论文研究了什么问题，更要告诉读者这个研究得出的结论。例如：“在事实与虚构之间：梅乐、卡彭特、沃尔夫的新闻观”就比“三个美国作家的新闻观研究”更专业、更准确。

\section{摘要的写法}
论文摘要是对论文研究内容的高度概括，应具有独立性和自含性，即应是一篇简短但意义完整的文章。通过阅读论文摘要，读者应该能够对论文的研究方法及结论有一个整体性的了解，因此摘要的写法应力求精确简明。论文摘要应包括对问题及研究目的的描述、对使用的方法和研究过程进行的简要介绍、对研究结论的高度凝练等，重点是结果和结论。

论文摘要切忌写成全文的提纲，尤其要避免“第 1 章……；第 2 章……；……”这样的陈述方式。

\section{引言的写法}
一篇学位论文的引言大致包含如下几个部分：1.~问题的提出；2.~选题背景及意义；3.~文献综述；4.~研究方法；5.~论文结构安排。
\begin{itemize}
    \item 问题的提出：要清晰地阐述所要研究的问题“是什么”。\footnote{选题时切记要有“问题意识”，不要选不是问题的问题来研究。}
    \item 选题背景及意义：论述清楚为什么选择这个题目来研究，即阐述该研究对学科发展的贡献、对国计民生的理论与现实意义等。
    \item 文献综述：对本研究主题范围内的文献进行详尽的综合述评，“述”的同时一定要有“评”，指出现有研究状态，仍存在哪些尚待解决的问题，讲出自己的研究有哪些探索性内容。
    \item 研究方法：讲清论文所使用的学术研究方法。
    \item 论文结构安排：介绍本论文的写作结构安排。
\end{itemize}

\section{正文的写法}
本部分是论文作者的研究内容，不能将他人研究成果不加区分地掺和进来。已经在引言的文献综述部分讲过的内容，这里不需要再重复。各章之间要存在有机联系，符合逻辑顺序。

\section{结论的写法}
结论是对论文主要研究结果、论点的提炼与概括，应精炼、准确、完整，使读者看后能全面了解论文的意义、目的和工作内容。结论是最终的、总体的结论，不是正文各章小结的简单重复。结论应包括论文的核心观点，主要阐述作者的创造性工作及所取得的研究成果在本领域中的地位、作用和意义，交代研究工作的局限，提出未来工作的意见或建议。同时，要严格区分自己取得的成果与指导教师及他人的学术成果。

在评价自己的研究工作成果时，要实事求是，除非有足够的证据表明自己的研究是“首次”“领先”“填补空白”的，否则应避免使用这些或类似词语。

\section{各节一级标题}
我是内容

\subsection{各节二级标题}
你是内容

\subsubsection{各节三级标题}
他是内容

\paragraph{四级标题}
内容内容

\subparagraph{五级标题}
内容内容

\section{字体样式}
宋体\quad \textbf{粗体}\quad \textit{斜体}\quad \textbf{\textit{粗斜体}}。

{\heiti 黑体\quad \textbf{粗体}\quad \textit{斜体}\quad \textbf{\textit{粗斜体}}}。

{\fangsong 仿宋\quad \textbf{粗体}\quad \textit{斜体}\quad \textbf{\textit{粗斜体}}}。

{\kaishu 楷书\quad \textbf{粗体}\quad \textit{斜体}\quad \textbf{\textit{粗斜体}}}。

Serif\quad \textit{Italic}\quad \textbf{Bold}\quad \textbf{\textit{BoldItalic}}

% {\sffamily Sans\quad \textit{Italic}\quad \textbf{Bold}\quad \textbf{\textit{BoldItalic}}}

{\ttfamily Mono\quad \textit{Italic}\quad \textbf{Bold}\quad \textbf{\textit{BoldItalic}}}
 % 引入第一章
\newpage % 另起一页

% 攻博期间发表的科研成果目录

在这里列出你攻博期间发表的科研成果 % 引入个人成果
\newpage % 另起一页

% 致谢页

以简短的文字表达作者对完成论文和学业提供帮助的老师、同学、领导、同事及亲属的感激之情。 % 引入致谢


\end{document}

